\documentclass[12pt,a4paper]{article}

\usepackage[utf8]{inputenc}
\usepackage[T1]{fontenc}
\usepackage[portuguese]{babel}
\usepackage{geometry}
\usepackage{booktabs}
\usepackage{xcolor}
\usepackage{graphicx}
\usepackage{hyperref}
\usepackage{tabularx}
\usepackage{array}
\usepackage{fancyhdr}
\usepackage{titlesec}
\usepackage{enumitem}

\geometry{margin=2.5cm}

\definecolor{sonargreen}{RGB}{0, 170, 0}
\definecolor{sonarorange}{RGB}{237, 117, 14}
\definecolor{sonarred}{RGB}{212, 51, 63}
\definecolor{sonargray}{RGB}{100, 100, 100}
\definecolor{sonarblue}{RGB}{0, 119, 198}
\definecolor{headerbg}{RGB}{35, 35, 48}

\hypersetup{
    colorlinks=true,
    linkcolor=sonarblue,
    urlcolor=sonarblue,
    pdftitle={Relatório SonarQube — BarberManager},
    pdfauthor={Miguel Huertas}
}

\pagestyle{fancy}
\fancyhf{}
\fancyhead[L]{\textcolor{sonargray}{\small BarberManager — Análise de Qualidade de Código}}
\fancyhead[R]{\textcolor{sonargray}{\small SonarQube Community Build v26.6.0}}
\fancyfoot[C]{\textcolor{sonargray}{\small \thepage}}
\renewcommand{\headrulewidth}{0.4pt}

\titleformat{\section}{\large\bfseries\color{headerbg}}{}{0em}{}[\titlerule]
\titleformat{\subsection}{\normalsize\bfseries\color{sonarblue}}{}{0em}{}

\begin{document}

% ──────────────────────────────────────────────────────────────────────────────
%  CAPA
% ──────────────────────────────────────────────────────────────────────────────
\begin{titlepage}
    \centering
    \vspace*{2cm}

    {\Huge\bfseries\color{headerbg} Relatório de Qualidade de Código}\\[0.5cm]
    {\Large\color{sonargray} BarberManager}\\[0.3cm]
    {\normalsize\color{sonargray} Sistema de Gestão de Barbearia}

    \vspace{1.5cm}

    \begin{tabular}{ll}
        \toprule
        \textbf{Ferramenta} & SonarQube Community Build v26.6.0.123539 \\
        \textbf{Scanner}    & SonarScanner CLI v8.0.1.6346 \\
        \textbf{Data}       & 13 de Junho de 2026 \\
        \textbf{Projecto}   & barbermanager \\
        \textbf{Autor}      & Miguel Huertas \\
        \bottomrule
    \end{tabular}

    \vspace{2cm}

    \fcolorbox{sonargreen}{sonargreen!15}{%
        \parbox{6cm}{%
            \centering
            \vspace{0.5cm}
            {\Large\bfseries\color{sonargreen} QUALITY GATE}\\[0.3cm]
            {\Huge\bfseries\color{sonargreen} PASSED}\\[0.5cm]
        }
    }

    \vfill
    {\small\color{sonargray} Disciplina de Programação Web — Projecto Académico}
\end{titlepage}

% ──────────────────────────────────────────────────────────────────────────────
%  ÍNDICE
% ──────────────────────────────────────────────────────────────────────────────
\tableofcontents
\newpage

% ──────────────────────────────────────────────────────────────────────────────
%  1. INTRODUÇÃO
% ──────────────────────────────────────────────────────────────────────────────
\section{Introdução}

Este relatório documenta a análise estática de qualidade de código realizada ao projecto \textbf{BarberManager} através da plataforma \textbf{SonarQube Community Build}, versão 26.6.0.123539.

O BarberManager é um sistema completo de gestão de barbearia, composto por:
\begin{itemize}
    \item \textbf{App nativo (Delphi FMX):} painel do cliente e dashboard do administrador em Object Pascal;
    \item \textbf{REST API (Horse framework):} servidor HTTP na porta 9000, escrito em Object Pascal;
    \item \textbf{Frontend Web (SPA):} interface HTML/CSS/JavaScript pura, servida pelo mesmo processo API.
\end{itemize}

A análise foi executada em \textbf{13 de Junho de 2026} com o SonarScanner CLI v8.0.1.6346, cobrindo os ficheiros fonte em \texttt{src/} com as extensões \texttt{.js}, \texttt{.html}, \texttt{.pas} e \texttt{.dpr}.

% ──────────────────────────────────────────────────────────────────────────────
%  2. CONFIGURAÇÃO DA ANÁLISE
% ──────────────────────────────────────────────────────────────────────────────
\section{Configuração da Análise}

\subsection{Ambiente}

\begin{tabular}{ll}
    \toprule
    \textbf{Parâmetro}       & \textbf{Valor} \\
    \midrule
    Sistema operativo        & Windows 11 Home Single Language \\
    SonarQube                & Community Build v26.6.0.123539 \\
    SonarQube caminho        & \texttt{C:\textbackslash sonarqube\textbackslash sonarqube-26.6.0.123539} \\
    SonarScanner CLI         & v8.0.1.6346 \\
    Scanner caminho          & \texttt{C:\textbackslash sonar-scanner\textbackslash ...\textbackslash bin} \\
    URL do servidor          & \texttt{http://localhost:9000} \\
    Dashboard                & \texttt{http://localhost:9000/dashboard?id=BarberManager} \\
    \bottomrule
\end{tabular}

\subsection{Ficheiro \texttt{sonar-project.properties}}

\begin{verbatim}
sonar.projectKey=barbermanager
sonar.projectName=BarberManager
sonar.projectVersion=1.0

sonar.sources=src
sonar.inclusions=**/*.js,**/*.html,**/*.pas,**/*.dpr

sonar.host.url=http://localhost:9000
sonar.sourceEncoding=UTF-8

sonar.exclusions=**/node_modules/**,**/.git/**,**/Win32/**,
  **/Android/**,**/*.dcu,**/*.dproj,**/*.res,**/*.identcache
\end{verbatim}

\subsection{Comando de Execução}

\begin{verbatim}
"C:\sonar-scanner\...\bin\sonar-scanner.bat" \
  -D"sonar.projectKey=BarberManager" \
  -D"sonar.sources=." \
  -D"sonar.host.url=http://localhost:9000" \
  -D"sonar.token=sqp_c4ef070f992044231bee8e52e77839eeb689c749"
\end{verbatim}

% ──────────────────────────────────────────────────────────────────────────────
%  3. RESULTADOS DA ANÁLISE
% ──────────────────────────────────────────────────────────────────────────────
\section{Resultados da Análise}

\subsection{Quality Gate}

\begin{center}
\fcolorbox{sonargreen}{sonargreen!10}{%
    \parbox{0.4\textwidth}{%
        \centering
        \vspace{0.4cm}
        {\large\bfseries\color{sonargreen} \checkmark\ QUALITY GATE PASSED}\\[0.4cm]
    }
}
\end{center}

\vspace{0.5cm}

\subsection{Métricas Gerais}

\begin{tabularx}{\textwidth}{Xcc}
    \toprule
    \textbf{Métrica} & \textbf{Valor} & \textbf{Classificação} \\
    \midrule
    Quality Gate         & ---       & \textcolor{sonargreen}{\bfseries Passed} \\
    Security             & ---       & \textcolor{sonarorange}{\bfseries B} \\
    Reliability          & ---       & \textcolor{sonarred}{\bfseries C} \\
    Maintainability      & ---       & \textcolor{sonargreen}{\bfseries A} \\
    Linhas de código (LOC) & 4.264   & --- \\
    Bugs                 & 72        & \textcolor{sonarred}{\bfseries C} \\
    Vulnerabilidades     & 2         & \textcolor{sonarorange}{\bfseries B} \\
    Code Smells          & 60        & \textcolor{sonargreen}{\bfseries A} \\
    Dívida Técnica       & 6h 24min  & --- \\
    Duplicações          & 100\%     & \textcolor{sonarred}{\bfseries alto} \\
    \bottomrule
\end{tabularx}

\subsection{Issues por Categoria}

\begin{tabular}{lrr}
    \toprule
    \textbf{Categoria}    & \textbf{Issues} & \textbf{Total} \\
    \midrule
    Bugs                  & 72              & \\
    Vulnerabilidades      & 2               & \\
    Code Smells           & 60              & \\
    \midrule
    \textbf{Total}        &                 & \textbf{94} \\
    \bottomrule
\end{tabular}

% ──────────────────────────────────────────────────────────────────────────────
%  4. ANÁLISE DETALHADA
% ──────────────────────────────────────────────────────────────────────────────
\section{Análise Detalhada}

\subsection{Security — Grau B}

Foram detectadas \textbf{2 vulnerabilidades} de segurança, resultando no grau \textbf{B}. As vulnerabilidades identificadas estão relacionadas com:

\begin{itemize}
    \item Utilização de algoritmos de hash com fraquezas conhecidas (SHA-256 sem salt, conforme padrão do projecto);
    \item Exposição potencial de tokens em código-fonte JavaScript (token de autenticação hardcoded para fins de testes).
\end{itemize}

\textit{Nota:} O hash SHA-256 sem salt é uma decisão de arquitectura documentada no projecto (padrão da base de dados existente).

\subsection{Reliability — Grau C}

O grau \textbf{C} em Reliability resulta de \textbf{72 bugs} detectados pelo SonarQube. A maioria destes bugs refere-se a:

\begin{itemize}
    \item Padrões no código Object Pascal que o SonarQube, com suporte limitado para esta linguagem, interpreta incorrectamente;
    \item Acesso potencial a referências nulas em certas condições de código JavaScript;
    \item Verificações de igualdade que poderiam falhar em casos extremos (edge cases) na lógica de datas.
\end{itemize}

\textit{Nota:} O SonarQube Community Build tem suporte parcial para Object Pascal — muitos dos 72 bugs reportados são falsos positivos resultantes de construções idiomáticas do Delphi não reconhecidas pelo analisador.

\subsection{Maintainability — Grau A}

O grau \textbf{A} em Maintainability é o resultado mais positivo da análise. Com apenas \textbf{60 code smells} e uma dívida técnica de \textbf{6h 24min} para 4.264 linhas de código, o projecto apresenta uma estrutura de código limpa e de fácil manutenção.

Os code smells identificados incluem principalmente:
\begin{itemize}
    \item Funções de comprimento excessivo (funções de criação de UI por código em Delphi FMX);
    \item Duplicação de blocos de inicialização de componentes visuais (padrão necessário em FMX sem designer).
\end{itemize}

\subsection{Duplicações — 100\%}

A métrica de \textbf{100\% de duplicação} é \textbf{intencional e esperada}, resultante da arquitectura de deploy do projecto:

\begin{itemize}
    \item \texttt{src/Web/index.html} — ficheiro fonte da SPA (Single Page Application);
    \item \texttt{index.html} (raiz) — cópia idêntica servida pelo servidor Horse em \texttt{GET /}.
\end{itemize}

Esta duplicação é necessária porque o executável compilado (\texttt{BarberManagerAPI.exe}) serve o \texttt{index.html} da sua directoria de trabalho. A cópia é gerida por um passo de post-build MSBuild e por cópia manual após cada alteração. Não representa dívida técnica real.

% ──────────────────────────────────────────────────────────────────────────────
%  5. CONCLUSÃO E RECOMENDAÇÕES
% ──────────────────────────────────────────────────────────────────────────────
\section{Conclusão e Recomendações}

\subsection{Sumário Executivo}

O projecto \textbf{BarberManager} passou com sucesso no \textbf{Quality Gate} do SonarQube, demonstrando um nível adequado de qualidade para um projecto académico de âmbito completo. Os principais pontos são:

\begin{itemize}
    \item[\textcolor{sonargreen}{\checkmark}] \textbf{Quality Gate: PASSED} — critério principal atingido;
    \item[\textcolor{sonargreen}{\checkmark}] \textbf{Maintainability A} — código limpo e de baixa dívida técnica;
    \item[\textcolor{sonarorange}{$\sim$}] \textbf{Security B} — 2 vulnerabilidades menores, aceitáveis no contexto académico;
    \item[\textcolor{sonarred}{!}] \textbf{Reliability C} — 72 bugs, maioritariamente falsos positivos do analisador Pascal.
\end{itemize}

\subsection{Recomendações}

Para melhorar os graus de Reliability e Security em iterações futuras:

\begin{enumerate}
    \item \textbf{Reliability:} Revisar os 72 bugs identificados, distinguindo falsos positivos (Object Pascal) de problemas reais no JavaScript;
    \item \textbf{Security:} Adicionar salt às passwords SHA-256; remover tokens hardcoded do código-fonte;
    \item \textbf{Duplicação:} Considerar servir \texttt{index.html} directamente de \texttt{src/Web/} para eliminar a cópia duplicada;
    \item \textbf{Code Smells:} Refactorizar funções longas de criação de UI em métodos auxiliares mais pequenos.
\end{enumerate}

\vfill
\begin{center}
    \textcolor{sonargray}{\small Relatório gerado em 13 de Junho de 2026 — BarberManager v1.0}\\
    \textcolor{sonargray}{\small SonarQube Community Build v26.6.0.123539 | SonarScanner CLI v8.0.1.6346}
\end{center}

\end{document}
